% Detailed methods moved from the main Methods.
\subsection*{Sequence encoder and supervised training}
Our sequence models build on CORE-BEHRT \citep{odgaard2024core}, adapted to a
causal encoder with next-event pretraining and multi-outcome competing-risk
prediction. Each patient's history was represented as a chronological sequence
of diagnoses, procedures, medications, laboratory results (discretised into
value bins) and hospital contacts, preceded by sex and age. A causal transformer
encoder \citep{vaswani2017attention} with a hidden size of 256, four layers and
four attention heads, and a scaled Time2Vec encoding of time
\citep{kazemi2019time2vec}, was pretrained to predict the next event over a
context of 8,192 events. \placeholder{Number of individuals and events,
training budget and selected epoch.} This encoder initialised every supervised
model in the main comparisons.

\subsection*{Supervised training across diseases and outcomes}
Every supervised model consisted of the pretrained encoder, trained further,
and a linear head that produced, for each outcome, piecewise-constant hazards
for the outcome and for competing death, building on neural interval-hazard and
competing-risk models \citep{gensheimer2019nnetsurvival,lee2018deephit}.
Cumulative incidence at each horizon was derived from these hazards. The loss
was a binary cross-entropy at each horizon: patients with the event before the
horizon were positive, patients who died of another cause or were followed
beyond the horizon were negative, and patients censored earlier did not
contribute at that horizon. Losses were averaged within each outcome and then
across endpoints, where the grade variants of one laboratory finding (grade
$\geq$2 and $\geq$3) and nested definitions of the same finding formed one
endpoint (69 outcomes, 52 endpoints), so that a finding defined at two grades
did not count twice. This weighting is separate from the clinical domains used
to present results, which group related outcomes for display only.

The four supervision scopes differed only in which cells of the
disease-by-outcome grid---one outcome in one disease---provided labels
(\cref{fig:supervision-grid}); we use ``cell'' for such a combination
throughout. A specialist learned one outcome
in one disease; a disease-pooled model learned one outcome across all diseases;
an outcome-pooled model learned all outcomes within one disease; and the joint
model learned all outcomes in all diseases. Every model predicted all six
horizons and was trained for up to 30 epochs, stopping after five epochs
without improvement in tuning loss, with the same optimiser, batch size and
sampling settings, and every comparison used the same target patients, labels
and horizons. Because equal epochs do not mean equal optimisation steps or
computation, we report both for each scope.

\subsection*{Outcome panel}
We predicted 69 outcomes covering disease course (second-line treatment,
treatment failure, transformation, new malignancy and death), infection,
haematological toxicity and transfusion, kidney, liver and coagulation
toxicity, electrolyte and metabolic disturbances, cardiovascular events and
acute care, each at 30, 60, 90, 180, 365 and 730 days. Laboratory outcomes used
grading thresholds aligned with, but not identical to, the Common Terminology
Criteria for Adverse Events \citep{nci2025ctcae}, with baseline-aware thresholds
when pre-treatment values were already abnormal; infection was additionally
captured through antimicrobial treatment, which we label as such. A patient
contributed to an outcome only if it was ascertainable for them---for example,
if laboratory follow-up was available---and was otherwise treated as missing
rather than negative. Death was a competing event for every non-fatal outcome.
The 69 outcomes were selected from 87 candidates by a prespecified rule based on
training and tuning event counts at one year, without held-out data
(Supplementary Methods S8). Definitions, code lists and risk-set sizes are
given in Supplementary Methods S5--S9.

\subsection*{Tabular comparators: features and tuning}
For each disease--outcome cell and horizon, we trained gradient-boosted trees
(XGBoost; \citealp{chen2016xgboost}) and ridge-penalised logistic regression on the
same patients, labels and competing-event handling as the neural models. Their
features came from the pipeline of our earlier DLBCL study
\citep{werling2026dlbcl}: using timeseriesflattener
\citep{bernstorff2023tsflattener}, prescriptions, laboratory results,
procedures and referrals, pathology codes, diagnoses, blood cultures and
microbiology were summarised over 90-day, one-year and three-year windows before
the index date, with aggregates suited to each source (counts and distinct
counts for codes; latest, minimum, maximum, variance and slope for laboratory
values). Variables recorded for fewer than 10\% of training patients were
represented by per-source counts of such rare occurrences. To give the tabular
and sequence models the same information, quality-registry variables such as
stage and performance status were removed, and age and sex at the index date
were added. Unlike the earlier study, no LASSO preselection was applied.
Numeric features were median-imputed with missingness indicators and
standardised, and categorical features one-hot encoded, with all preprocessing
fitted on training data. Hyperparameters were chosen for each cell by log loss
on the tuning set from a prespecified grid: for XGBoost, shallow trees (depth
1--3), deeper trees (depth 4 and 6) and the setting of the earlier study
(depth 8, learning rate 0.01, up to 3,000 trees), with the number of trees also
chosen on the tuning set; for logistic regression, the penalty strength along
a seven-value path ($C$ from $3\times10^{-4}$ to 1, fitted with warm starts). Neither model was class-reweighted, so both produce
probabilities that can be compared directly. \placeholder{Confirm against the
delivered feature export: anchoring at first-line treatment, source list and
number of features.}

As a reference outside the information-matched comparison, XGBoost was also
trained for each cell with the clinical quality-registry variables recorded at
diagnosis added to the same features (performance status, stage, lactate
dehydrogenase, extranodal involvement, B symptoms and cytogenetic risk where the
registry for the disease records them; \placeholder{final list}). Values were
taken as recorded at diagnosis, left missing where not recorded, and handled by
the same missingness indicators; hyperparameters were tuned on the same grid.
The sequence models never saw these variables, so this arm estimates what the
curated registry adds, not a like-for-like comparison.

\subsection*{Comparison with clinical prognostic indices}
The quality registries record, for each patient, the prognostic index scored at
diagnosis (for example the IPI, FLIPI, MIPI, CLL-IPI or ISS) and its risk
category (low, intermediate or high), before treatment started. Every index
recorded for at least 50 patients of a subtype, and every index used in more
than one subtype pooled across all subtypes recording it, was compared with the joint
model's predicted two-year risk of treatment failure, with death part of the
endpoint, in the patients who had that index. For DLBCL we also computed the
NCCN-IPI from the components recorded in the lymphoma registry, following
the original definition \citep{zhou2014nccnipi}: age (41--60, 61--75 and
over 75 years: 1, 2 and 3 points), LDH relative to the upper limit of normal
(205~U/L, or 255~U/L from age 70; ratio above 1 and above 3: 1 and 2 points),
Ann Arbor stage III--IV (1 point), involvement of bone marrow, central nervous
system, liver, gastrointestinal tract or lung (1 point), and performance
status $\geq$2 (1 point); low 0--1, low-intermediate 2--3, high-intermediate
4--5, high $\geq$6. Patients missing any component were not scored.

Model risk groups were fixed on tuning patients with known two-year status and
applied unchanged to the test set. The low cutoff was the highest predicted
risk at which the group below it, of at least 20 patients, had an observed risk
no higher than the index's low-risk category. The high cutoff was the predicted
risk at which the model's specificity equalled that of the index's high-risk
category, so that both flagged the same proportion of patients without the
event; this is also the operating point marked in the precision--recall
comparison. Patients between the two cutoffs were divided in the proportions
of the index's intermediate categories (one for the registry indices, two for
the NCCN-IPI). Anchoring both ends to the index makes the groups directly
comparable while leaving the model free to differ: a model that ranks patients
better places more patients at low risk at the same risk, and captures more
events at the same specificity. For each index and each system we report group
size, observed risk, and the sensitivity and specificity of classifying a
patient in that group or higher, and, for the model, the observed risk in its
highest-risk 5\% and 10\% of patients, which the index cannot resolve within
its top category. As a sensitivity analysis, model groups were instead given
exactly the index's group sizes. Discrimination used the index category (or
NCCN-IPI points) as an ordinal score.

\subsection*{Routing formulation and selection}
To test whether a joint model benefits from conditioning its predictions on
disease and outcome beyond ordinary shared learning, we compared the joint
model with a routed model trained on the same panel. In the routed model, the
patient representation
$h_i$ is adjusted for each target outcome $k$ by a weighted combination of
small residual adapters, one per disease and one per outcome:
\begin{equation}
 \begin{aligned}
 h_{ik}&=h_i+\lambda\sum_s\pi_{iks}\,c_s\,A_s(h_i),\\
 A_s(h)&=U_s\,\mathrm{GELU}(V_s h),
 \qquad c_s=\frac{e_s}{e_s+\tau}.
 \end{aligned}
 \label{eq:router}
\end{equation}
where $A_s$ is a low-rank adapter for source $s$, the weights $\pi_{iks}$ are
computed from the patient, the patient's disease and the target outcome, and
$c_s$ shrinks the contribution of sources with few training events $e_s$. The
adapted representation enters the same linear head as in the joint model. Each adapter was
updated only by the supervision of its own disease or outcome. The approach
relates to mixtures of experts and parameter-efficient adapters
\citep{jacobs1991adaptive,ma2018mmoe,houlsby2019adapters,pfeiffer2021adapterfusion}.

The router configuration was chosen on the tuning set in a short prespecified
sequence: which adapter banks to include (none, disease, outcome or both), then
the adapter rank, then the form of the routing query. A simpler configuration
was preferred unless a more complex one improved tuning performance beyond a
prespecified margin, and a routed model with no adapters was required to
match the joint model before differences between adapter banks were interpreted.
\placeholder{Insert the selected configuration.} Because a target outcome
already has its own prediction head, an outcome adapter bank must add value
beyond that head; route weights describe how the model allocates corrections
and are not interpreted as evidence of relationships between diseases.



\begin{figure}[!htbp]
  \centering
  \resizebox{0.96\textwidth}{!}{\begin{tikzpicture}[opera,x=1cm,y=1cm]
  % One target prediction, represented as four clearly separated stages.
  \node[eyebrow,anchor=west] at (0,0) {SOURCE-ANCHORED ROUTER};
  \node[labelsmall,anchor=west] at (0,-.36)
    {One patient $i$ and one target outcome $k$};

  \node[eyebrow] at (1.35,-1.02) {INPUTS};
  \node[eyebrow] at (5.02,-1.02) {ROUTE};
  \node[eyebrow] at (9.28,-1.02) {SOURCE BANK};
  \node[eyebrow] at (14.10,-1.02) {PREDICTION};
  \draw[divider] (3.03,-.88)--(3.03,-5.08);
  \draw[divider] (7.02,-.88)--(7.02,-5.08);
  \draw[divider] (11.94,-.88)--(11.94,-5.08);

  % Input identities are not source representations: the encoder supplies z.
  \node[stageblue,minimum width=2.18cm] (patient) at (1.38,-2.08)
    {patient state $z_i$};
  \node[stageteal,minimum width=2.18cm] (disease) at (1.38,-3.08)
    {fine disease $d_i$};
  \node[stagewarm,minimum width=2.18cm] (target) at (1.38,-4.08)
    {target outcome $k$};

  \node[stageblue,minimum width=2.48cm,minimum height=1.04cm]
    (router) at (5.03,-3.08) {learned weights\\$\pi_{i,k,s}$};
  \draw[flowblue] (patient.east) to[out=0,in=160] (router.north west);
  \draw[flowteal] (disease.east)--(router.west);
  \draw[flowwarm] (target.east) to[out=0,in=200] (router.south west);
  \node[captionnote,text width=3.40cm] at (5.03,-4.36)
    {Patient, disease and target context};

  % The three rows are categories, not a claim that only three sources exist.
  \node[stageblue,minimum width=3.58cm] (shared) at (9.45,-1.97)
    {$z_i$ \quad shared path};
  \node[stageteal,minimum width=3.58cm] (dsrc) at (9.45,-3.08)
    {$\sum_d\pi_{i,k,d}\,a_d(z_i)$ \quad disease};
  \node[stagewarm,minimum width=3.58cm] (osrc) at (9.45,-4.19)
    {$\sum_o\pi_{i,k,o}\,a_o(z_i)$ \quad outcome};
  \draw[flowblue] (router.east) to[out=25,in=180] (shared.west);
  \draw[flowteal] (router.east)--(dsrc.west);
  \draw[flowwarm] (router.east) to[out=-25,in=180] (osrc.west);

  \node[stage,minimum width=2.33cm] (mix) at (14.07,-2.56)
    {weighted residual\\$z_{i,k}$};
  \node[outcome,minimum width=2.33cm] (hazards) at (14.07,-4.02)
    {interval hazards\\$\lambda_{i,k,t}$};
  \draw[flowblue] (shared.east) to[out=0,in=155] (mix.north west);
  \draw[flowteal] (dsrc.east)--(mix.west);
  \draw[flowwarm] (osrc.east) to[out=0,in=205] (mix.south west);
  \draw[flowwarm] (mix.south)--(hazards.north);

  \draw[divider] (0,-5.34)--(15.84,-5.34);
  \node[captionnote,anchor=west,text width=15.40cm,align=left]
    at (.20,-5.68)
    {The shared path remains unscaled; route weights govern adapter residuals. The bank includes eligible borrowed sources, whose adapter parameters are detached.};
\end{tikzpicture}
}
  \caption{\textbf{Source-anchored disease--outcome routing.} For each patient
  and target outcome, the router mixes a shared representation with residual
  adapters indexed by fine disease and outcome. Its weights depend on the
  patient, disease and target; sources with no training support are masked and
  low-event sources are shrunk. An adapter is updated by its own eligible
  supervision, while borrowed adapter parameters are detached. The mixed
  representation enters a shared linear head for interval hazards and
  cumulative incidence. This is a model schematic, not an empirical route-weight
  or performance plot.}
  \label{fig:source-router}
\end{figure}


\subsection*{Hierarchy for effect heterogeneity}
The primary evidence is the paired effect in each disease--outcome cell with
its patient-bootstrap interval. Many cells are small, so their raw estimates
are noisy and rankings of individual cells are unreliable. We therefore also
fit a deliberately simple Bayesian hierarchical model whose purpose is partial
pooling: each cell's effect is the panel mean plus a disease effect, an outcome
effect and a cell-specific deviation, so that a noisy cell is shrunk towards
what similar cells show, by an amount the data determine. The model answers
two questions: how large is the typical effect, and how much does it vary
between diseases and between outcomes? The cell estimates are not treated as
independent: their sampling covariance, taken from the shared patient-bootstrap
draws, enters the likelihood, because the same patients contribute to every
outcome within a disease.

Whether effects depend on the amount of training data available for the
target is a separate question; we do not assume a linear relationship.
A low-dimensional spline still constrains smoothness and complexity. The cell effect is allowed to vary smoothly with log training
events through a cubic spline with few knots and a shrinkage prior, to regularise departures from a flat relationship; the prior still affects
the fitted shape and is assessed in sensitivity checks. As a transparent check
that does not depend on smoothing, cells are also grouped into prespecified
bins of training events (fixed before held-out evaluation) and the mean effect
is estimated within each bin. Within each
question, the model fits the independent effects directly (for pooling:
disease pooling, outcome pooling and their interaction); conditional and
combined contrasts are exact sums of these, so all reported pooling contrasts
are mutually consistent. Priors
are weakly informative on each metric's scale and checked with prior
predictive simulation; convergence and posterior predictive checks are
reported. Posterior probabilities of benefit, harm and practical equivalence
use a prespecified tolerance. Raw paired estimates are always shown alongside.
\placeholder{Insert locked priors, tolerances, support-bin edges and sampling
diagnostics.}


\FloatBarrier
